\documentclass[10pt,a4paper]{article}
\usepackage{fontspec}
\setmainfont{texgyretermes}[Extension=.otf,UprightFont=*-regular,BoldFont=*-bold,ItalicFont=*-italic,BoldItalicFont=*-bolditalic]
\usepackage[margin=1.8cm,headheight=14pt]{geometry}
\usepackage{amsmath,amssymb,booktabs,tabularx,array,ragged2e,enumitem,microtype,xcolor,tikz,fancyhdr,titlesec}
\usetikzlibrary{arrows.meta,positioning,fit,backgrounds}
\usepackage[hidelinks,unicode]{hyperref}
\definecolor{ink}{HTML}{183A45}\definecolor{teal}{HTML}{087F7A}
\definecolor{blue}{HTML}{405F8A}\definecolor{gray}{HTML}{56656B}
\newcolumntype{P}[1]{>{\RaggedRight\arraybackslash}p{#1}}
\newcolumntype{Y}{>{\RaggedRight\arraybackslash}X}
\renewcommand{\tabularxcolumn}[1]{>{\RaggedRight\arraybackslash}p{#1}}
\setlength{\parindent}{0pt}\setlength{\parskip}{5pt}\setlength{\emergencystretch}{2em}
\setlength{\tabcolsep}{4pt}\renewcommand{\arraystretch}{1.12}
\setlist{nosep,leftmargin=1.4em,topsep=3pt}
\titleformat{\section}{\large\bfseries\color{ink}}{\thesection}{.6em}{}
\titleformat{\subsection}{\normalsize\bfseries\color{ink}}{\thesubsection}{.6em}{}
\titlespacing*{\section}{0pt}{8pt}{6pt}\titlespacing*{\subsection}{0pt}{7pt}{4pt}
\pagestyle{fancy}\fancyhf{}
\fancyhead[L]{\small\color{gray}GEOI-BQ: PHƯƠNG PHÁP, BẢO ĐẢM VÀ THỰC NGHIỆM}
\fancyhead[R]{\small\color{gray}27/10/2026}\fancyfoot[C]{\small\thepage}
\renewcommand{\headrulewidth}{.3pt}
\newcommand{\key}[1]{\par\smallskip\noindent\textcolor{teal}{\textbf{#1}}\par\smallskip}
\newcommand{\source}[1]{\par{\footnotesize\color{gray}Nguồn: #1}\par}
\newcommand{\claim}[1]{\par\textbf{Kết luận cần giữ.} #1\par}
\newcommand{\TV}{\operatorname{TV}}
\hypersetup{pdftitle={GeoI-BQ: phương pháp, bảo đảm riêng tư và đánh giá thực nghiệm},pdfauthor={}}
\newcommand{\RatioUpper}{2{,}437372}
\newcommand{\PosteriorUpper}{70{,}91}
\newcommand{\PrimaryRecall}{99{,}595}

\begin{document}
\begin{center}
{\LARGE\bfseries\color{ink}Bảo vệ quỹ đạo trong dịch vụ vị trí}\\[3pt]
{\large GeoI-BQ: từ đặc tả thuật toán đến bảo đảm toàn transcript}
\end{center}

% PAGE 01
\section{Bài toán, mục tiêu và phạm vi}
\textbf{Bài toán.} Người dùng di chuyển và tìm POI qua một nhà cung cấp LBS. Dù không nhận GPS thật,
nhà cung cấp vẫn có thể ghép các yêu cầu theo thời gian để suy luận vị trí và đường đi.
Mục tiêu là giới hạn thông tin bổ sung từ \emph{toàn transcript}, đồng thời giữ chất lượng đáp án và đo chi phí phục vụ.

\textbf{Mạch lập luận.} Đặc tả attacker và output trước; mô tả cơ chế chạy; phát biểu giả thiết và chứng minh;
sau đó dùng benchmark để đo đánh đổi. Các vòng phát triển cũ không được ghép thành một cấu hình duy nhất.

\subsection{Nền tảng và khoảng trống cần xử lý}
Geo-indistinguishability (Geo-I) là thuộc tính ràng buộc tỷ số xác suất đầu ra theo khoảng cách GPS.
REM lấy mẫu trực tiếp trên miền đường công khai để đáp ứng thuộc tính đó trong mô hình lý tưởng.
Với quỹ đạo, cần tính cả phép kiểm tra giữ/đổi Z, tổng chi phí nhiều bước và cách thông tin đã bảo vệ
được chuyển thành các Q. Chọn Q ngẫu nhiên có giới hạn bổ sung một cận toàn phiên riêng.

Geo-I kế thừa \href{https://arxiv.org/abs/1212.1984}{Andrés và cộng sự (CCS 2013)};
private test và predictive budget có tiền lệ ở
\href{https://petsymposium.org/2014/papers/Chatzikokolakis.pdf}{Chatzikokolakis và cộng sự (PETS 2014)}.
Không coi hợp thành cho trajectory hoặc noisy reuse tự nó là đóng góp mới của luận văn.

Các paper có thể trả một tọa độ, một quỹ đạo ảo, hoặc một tập dummy có/không chứa vị trí thật.
Vì thế entropy, xác suất chọn điểm thật và sai số tái dựng không tự tương đương. Luận văn dùng metric
theo đúng tác vụ, giữ metric gốc khi đủ dữ liệu và không ép mọi dạng output vào một score.
Các adaptation lịch sử và metric khả dụng được giữ riêng ở Phụ lục A.

\subsection{Attacker và ranh giới quan sát}
Thiết bị giữ GPS, Z, belief, yêu cầu thật và randomness. Server biết thuật toán, cấu hình, bản đồ,
đồng hồ công khai; thấy Q có thứ tự, request chung, phản hồi và retry logic.
Clock mở/đóng phiên được giữ chung trong phép so sánh. Các bảo đảm không bao gồm kênh account/IP/biển số,
hoặc thời điểm gửi/lỗi phụ thuộc GPS ngoài giao thức đã đặc tả.
Traces được cố định trước phép chạy; trường hợp đường đi thật thay đổi do POI gợi ý cần mô hình tương tác riêng.

\begin{tabularx}{\linewidth}{@{}P{2.4cm}P{4.2cm}Y@{}}\toprule
\textbf{Scenario} & \textbf{Mục tiêu suy luận} & \textbf{Vai trò trong bản này}\\\midrule
S1, S2, S3 & Vị trí hiện tại, nơi dừng, đường đã đi & Benchmark vị trí; S2 mới là proxy tọa độ cửa sổ muộn, chưa là phát hiện nơi dừng ngữ nghĩa.\\
S9, S10 & Điểm đầu và điểm cuối & Benchmark endpoint dưới cùng protocol; không suy ẩn thời điểm bắt đầu/kết thúc.\\
S4, S5, S6 & Liên kết người/xe, cạnh tiếp theo, đích tương lai & Hệ quả suy luận có điều kiện, tách khỏi kết quả thực nghiệm chính.\\
S7 & Nội dung và ý định truy vấn & Request chung; nhu cầu giữ local. Không đồng nhất với che mọi ý định tương quan tuyến.\\
S8 & Thông tin người đồng hành & Mở rộng; chưa tuyên bố bảo vệ nhóm cho cấu hình hiện tại.\\\bottomrule
\end{tabularx}

\textbf{Câu hỏi nghiên cứu.} (RQ1) Chứng minh được cận privacy nào cho toàn transcript?
(RQ2) Utility được bảo đảm ở mức nào và với điều kiện gì? (RQ3) Các đánh đổi privacy--utility--cost
thể hiện ra sao trước baseline, control và ablation?
\claim{Không gửi GPS hoặc có năm Q giả chưa đủ để kết luận bảo vệ quỹ đạo. Cần cận cho phân phối toàn bộ đầu ra.}
\source{Brief hướng viết luận văn; đặc tả và cơ sở toán học GeoI-BQ v1.}

% PAGE 02
\clearpage
\section{Phương pháp GeoI-BQ}
\subsection{Kiến trúc, luồng dữ liệu và cấu hình chính}
\begin{center}
\begin{tikzpicture}[>=Latex,
 box/.style={draw=ink!60,fill=white,rounded corners=2pt,text width=8cm,align=center,inner sep=5pt,font=\small},
 ext/.style={draw=blue!65,fill=blue!4,rounded corners=2pt,text width=3cm,align=center,inner sep=5pt,font=\small},
 flow/.style={->,draw=ink,line width=.8pt},pub/.style={->,draw=gray,dashed},net/.style={->,draw=blue,line width=.8pt}]
\node[ext,text width=7.8cm] (gps) at (0,.7) {\textbf{ĐẦU VÀO: GPS thật $x_t$}\\Giữ tại thiết bị};
\node[box] (gi) at (0,-.9) {\textbf{Lớp 1. Tạo và duy trì thông tin đã bảo vệ}\\REM tạo Z trên đường; noisy test giữ hoặc đổi Z\\Private read cách ít nhất 60 s; cap Geo-I từng phiên};
\node[box] (belief) at (0,-2.5) {\textbf{Lớp 2. Ước lượng belief $b_t$}\\Lịch sử Z + thời gian + chuyển động gần đúng\\Trọng số trên các ô/vị trí đại diện};
\node[box,draw=teal,line width=1pt] (q) at (0,-4.1) {\textbf{Lớp 3. Chọn một đoạn Q theo xác suất}\\Thư viện khả thi và xác suất nền từ thông tin công khai\\Belief chỉ điều chỉnh xác suất trong giới hạn};
\node[box] (send) at (0,-5.7) {\textbf{Lớp 4. Phát các khung đã cam kết}\\Ba khung 0/20/40 s; mỗi khung có $K=5$ Q\\Truy vấn mọi loại POI, $L=60$ mỗi loại và mỗi Q};
\node[box] (local) at (0,-7.2) {\textbf{Lớp 5. Xử lý đáp án tại thiết bị}\\Union/cache hợp lệ; lọc và sắp xếp theo GPS + nhu cầu $\psi$};
\node[ext,text width=7.8cm,draw=teal] (out) at (0,-8.7) {\textbf{ĐẦU RA LOCAL: danh sách phù hợp đã sắp xếp}};
\node[ext] (context) at (6.4,-3.1) {\textbf{Công khai}\\Bản đồ, cấu hình, clock\\Q và phản hồi trước};
\node[ext] (server) at (6.4,-5.7) {\textbf{Server POI}\\Nhận Q, trả top-L\\Ngoài thiết bị};
\node[ext] (request) at (6.4,-8.0) {\textbf{Nhu cầu $\psi$}\\Tiêu chí, loại, bán kính, đích\\Giữ local};
\draw[flow] (gps)--(gi);\draw[flow] (gi)--(belief);\draw[flow] (belief)--(q);
\draw[flow] (q)--(send);\draw[flow] (send)--(local);\draw[flow] (local)--(out);
\draw[net] (send.east)--(server.west);\draw[net] (server.south)--(6.4,-6.6)--(local.east);
\draw[pub] (context.west)--(4.5,-3.1)--(4.5,-4.1)--(q.east);
\draw[pub] (local.west)--(-4.7,-7.2)--(-4.7,-4.1)--(q.west);
\draw[flow,draw=teal] (gps.west)--(-5.4,.7)--(-5.4,-7.2)--node[above,font=\footnotesize]{GPS local}(local.west);
\draw[flow] (request.west)--(local.south east);
\begin{scope}[on background layer]
\node[fit=(gi)(belief)(q)(send)(local),draw=teal,fill=teal!2,rounded corners=4pt,inner xsep=9pt,inner ysep=10pt] {};
\end{scope}
\node[anchor=south west,text=teal,fill=white,inner sep=1pt,font=\footnotesize\bfseries] at (-4.3,-.28) {MÔ HÌNH TẠI THIẾT BỊ};
\end{tikzpicture}
\end{center}
\begin{tabularx}{\linewidth}{@{}P{3.1cm}Y@{}}\toprule
\textbf{Thành phần} & \textbf{Cấu hình và ý nghĩa}\\\midrule
Geo-I & $B=0{,}24$ m$^{-1}$, $H=12$, $u=B/(2H)=0{,}01$ m$^{-1}$, $C_s=(2H-1)u=0{,}23$ m$^{-1}$.\\
Q và dịch vụ & $K=5$, $L=60$/loại/Q; gửi mỗi 20 s; chọn chung ba khung mỗi 60 s.\\
Ngân sách chọn đoạn & $\Gamma_Q=1$; $\varepsilon_{Q,j}=\Gamma_Q/[j(j+1)]$; ledger bền vững.\\
Planner công khai & Portfolio beam 24; novelty và long-range weight đều 4; thang travel 60 s.\\
Cache và ranking & Metadata tĩnh còn hợp lệ trong phiên; đáp án đầy đủ đã sắp xếp. Recall@5 chỉ dùng đánh giá.\\\bottomrule
\end{tabularx}
\claim{Z và belief không gửi đi. Request thật không điều khiển loại, độ sâu hoặc nhịp truy vấn server.}

% PAGE 03
\clearpage
\subsection{Từ GPS đến phân phối chọn đoạn Q}
Với $V$ là miền đường công khai và $a_v$ là tọa độ ứng viên, lần đầu lấy mẫu REM:
\begin{equation}
Z_t=a_{V_t},\quad V_t\sim R_u(\cdot\mid x_t),\qquad
R_u(v\mid x)=\frac{e^{-u\|x-a_v\|/2}}{\sum_{w\in V}e^{-u\|x-a_w\|/2}}.
\label{eq:rem}
\end{equation}
Sau đó, tại mốc được phép đọc GPS, cộng $\xi_t\sim\operatorname{Lap}(0,1/u)$ vào khoảng cách đến Z cũ.
Nếu $\|x_t-Z_{\rm cũ}\|+\xi_t>\theta$ thì tạo Z mới bằng REM; ngược lại giữ Z cũ. $\theta=200$ m.
Ngoài lịch hoặc không đủ dự toán, không đọc GPS mới cho lớp bảo vệ.

Belief dự đoán rồi cập nhật từ likelihood quan sát đã bảo vệ:
\begin{equation}
b_t^-(s)=\sum_r b_{t-1}(r)T_{\Delta t}(r,s),\qquad
b_t(s)=\frac{b_t^-(s)\ell_t(Z_t\mid s,Z_{\rm cũ})}{\sum_r b_t^-(r)\ell_t(Z_t\mid r,Z_{\rm cũ})}.
\label{eq:belief}
\end{equation}
Không có quan sát mới thì chỉ dự đoán. Likelihood tính cả noisy reuse và REM, có cộng bội số lane trùng tọa độ.
Belief là bộ lọc gần đúng, không phải posterior GPS đã được hiệu chuẩn.

Với mỗi đoạn khả thi $A_a=(Q_a^{(0)},Q_a^{(20)},Q_a^{(40)})$, tính điểm kỳ vọng và điểm cuối cùng:
\begin{equation}
v_a=\sum_s b(s)g_a(s),\qquad z_a=\frac{\operatorname{Quantize}(v)_a}{S}\in[0,1],\quad S=65536.
\label{eq:score}
\end{equation}
$g_a$ là bảng proxy công khai; raw scores được chuẩn hóa min--max trước lượng tử hóa.
Chữ $z_a$ thường là điểm số, khác tọa độ Z. Với xác suất nền công khai $w_a>0$, $\sum_aw_a=1$:
\begin{equation}
\mu=\sum_aw_az_a,\quad \alpha=\frac{\varepsilon_Q}{2+\varepsilon_Q},\qquad
p_a=w_a[1+\alpha(z_a-\mu)],\qquad A\sim\operatorname{Categorical}(p).
\label{eq:q}
\end{equation}
Thư viện và $w$ được cố định tại cùng lịch sử công khai; belief không được loại ứng viên. Đoạn có điểm
trên trung bình được tăng xác suất, sau đó cả đoạn được commit và phát theo clock.

\textbf{Ví dụ số minh họa.} Ba vùng có đại diện 0/100/300 m, prior $(0{,}5;0{,}3;0{,}2)$;
giả sử REM chỉ có ba đầu ra này và quan sát đầu tiên là $Z=100$ m, $u=0{,}01$:
\begin{small}\par\smallskip\noindent\begin{tabularx}{\linewidth}{@{}XXXX@{}}\toprule
\textbf{Vùng / đại diện} & \textbf{Prior} & \textbf{Likelihood Z=100} & \textbf{Belief sau} \\
\midrule
A / 0 m & 0{,}50 & 0{,}3315 & 45{,}54\% \\
B / 100 m & 0{,}30 & 0{,}5065 & 41{,}75\% \\
C / 300 m & 0{,}20 & 0{,}2312 & 12{,}71\% \\
\bottomrule\end{tabularx}\par\smallskip
\end{small}
Nếu hai đoạn có proxy $(0{,}9;0{,}6;0{,}2)$ và $(0{,}3;0{,}7;0{,}9)$ trên ba vùng,
điểm lần lượt khoảng $0{,}686$ và $0{,}543$. Chuẩn hóa cho $z=(1;0)$.
Với $w=(0{,}5;0{,}5)$ và $\varepsilon_Q=0{,}5$, ta có $\alpha=0{,}2$, $\mu=0{,}5$ và $p=(0{,}55;0{,}45)$.
Đây là tính toán minh họa, không phải một benchmark Recall.

\textbf{Ví dụ nhiều bước từ một TRAIN tape đã lưu.} Các thời điểm hiển thị thưa không bỏ qua những bước thực tế ở giữa:
\begin{small}\par\smallskip\noindent\begin{tabularx}{\linewidth}{@{}P{1.1cm}P{2.5cm}XXXX@{}}\toprule
\textbf{t (s)} & \textbf{GPS bảo vệ / nhánh} & \textbf{Chọn đoạn} & \textbf{ε Q mới} & \textbf{Q đã dành} & \textbf{POI cache} \\
\midrule
0 & Đọc / tạo mới & Có & 0{,}5000 & 0{,}5000 & 381 \\
20 & Không đọc & Không & -- & 0{,}5000 & 381 \\
40 & Không đọc & Không & -- & 0{,}5000 & 381 \\
60 & Đọc / tạo mới & Có & 0{,}1667 & 0{,}6667 & 385 \\
120 & Đọc / tạo mới & Có & 0{,}0833 & 0{,}7500 & 407 \\
300 & Đọc / giữ Z & Có & 0{,}0238 & 0{,}8571 & 418 \\
600 & Đọc / giữ Z & Có & 0{,}0076 & 0{,}9167 & 418 \\
\bottomrule\end{tabularx}\par\smallskip
\end{small}
Tại 20/40 s chỉ phát khung đã chọn; belief vẫn có thể dự đoán. Tại 300 s, test giữ Z nhưng vẫn tốn một đơn vị $u$.
Khi cache đã đủ 418 POI, dịch vụ gần hoàn chỉnh có thể che khuất lợi ích của thích nghi theo belief.

% PAGE 04
\clearpage
\section{Chứng minh privacy: từ thao tác đến toàn transcript}
\subsection{Mệnh đề 1: REM thỏa Geo-I}
Giữ cùng $V$ và cùng phép chiếu mét. Với hai GPS $x,x'$, đặt $d=\|x-x'\|$,
$w_x(v)=e^{-u\|x-a_v\|/2}$ và $S_x=\sum_vw_x(v)$.
Bất đẳng thức tam giác cho $|\|x-a_v\|-\|x'-a_v\||\le d$, suy ra
\begin{equation}
\frac{w_x(v)}{w_{x'}(v)}\le e^{ud/2},\qquad
w_{x'}(v)\le e^{ud/2}w_x(v)\ \Longrightarrow\ \frac{S_{x'}}{S_x}\le e^{ud/2}.
\end{equation}
Vì vậy, với mọi ứng viên $v$:
\begin{equation}
\boxed{\frac{R_u(v\mid x)}{R_u(v\mid x')}
=\frac{w_x(v)}{w_{x'}(v)}\frac{S_{x'}}{S_x}\le e^{ud}.}
\label{eq:remproof}
\end{equation}
Hoán đổi $x,x'$ cho cận ngược. Cộng theo các ứng viên trùng tọa độ giữ cận cho Z quan sát được.
Đây là tính chất của toàn kernel, không phải phép lọc riêng từng điểm ``đạt Geo-I''.

\subsection{Mệnh đề 2: noisy reuse và hạch toán toàn phiên}
Tại cùng lịch sử đã bảo vệ, Z cũ $h$ giống nhau. $g_h(x)=\|x-h\|$ là 1-Lipschitz.
Với $Y=g_h(x)+\operatorname{Lap}(0,1/u)$, mật độ là $f_x(y)=\frac u2e^{-u|y-g_h(x)|}$, nên
\begin{equation}
\frac{f_x(y)}{f_{x'}(y)}
=e^{u(|y-g_h(x')|-|y-g_h(x)|)}\le e^{ud}.
\label{eq:test}
\end{equation}
Tích phân trên $Y\le\theta$ và $Y>\theta$ chứng minh cả quyết định giữ và đổi thỏa cận $u$.
Gọi $r_x$ là xác suất giữ, kernel của lịch sử mở rộng có hai nhánh:
\[
K_x(\mathrm{reuse})=r_x,\qquad K_x(\mathrm{fresh},v)=(1-r_x)R_u(v\mid x).
\]
Nhánh reuse có hệ số tối đa $e^{ud}$; nhánh fresh có hệ số tối đa $e^{2ud}$.
Cờ nhánh ở local; phân tích lịch sử mạnh hơn rồi chiếu xuống output, không giả sử server nhìn thấy cờ này.

\begin{tabularx}{\linewidth}{@{}Ycc@{}}\toprule
\textbf{Thao tác} & \textbf{Đơn vị chi phí} & \textbf{Lý do}\\\midrule
Z đầu tiên bằng REM & 1 & Một lần REM\\
Test rồi giữ Z & 1 & Một lần test\\
Test rồi tạo mới & 2 & Test + REM\\
Chỉ hậu xử lý, không đọc GPS mới & 0 & Không thêm private input\\\bottomrule
\end{tabularx}

\textbf{Hợp thành.} So hai traces $X,X'$ trên cùng lịch \emph{có thể đọc}, kể cả mốc bị bỏ qua.
Prospective filter dự toán chi phí trước khi đọc; tại cùng lịch sử, pacing/cap quyết định giống nhau.
Với $c_i\in\{0,1,2\}$ và $\sum_ic_i\le M$, chain rule cho
\begin{equation}
\log\frac{\Pr(E=e\mid X,C)}{\Pr(E=e\mid X',C)}
\le u\sum_ic_i\|x_i-x_i'\|\le MuD_\infty=C_sD_\infty.
\label{eq:geosession}
\end{equation}
$D_\infty$ là độ lệch lớn nhất giữa hai traces ở các mốc tương ứng. Cận tinh hơn giữ tổng khoảng cách
của các thao tác khả thi; không cần độc lập giữa các GPS.

\textbf{Ý nghĩa ngân sách.} $u=B/(2H)$, $M=2H-1=23$, $C_s=Mu=0{,}23$ m$^{-1}$.
Đây là hệ quả cấu hình, không phải bán kính nhiễu hay mức an toàn tối ưu. Muốn tỷ số không vượt $R_*$
cho $D_\infty\le d_*$, điều kiện đủ của cận thô là $C_s\le\log(R_*)/d_*$.
Không khẳng định cấu hình hiện tại đã được chọn theo một chính sách rủi ro như vậy từ trước.
Hết cap vẫn có thể truy vấn bằng thông tin đã bảo vệ, nhưng khả năng thích nghi có thể giảm.

% PAGE 05
\clearpage
\subsection{Mệnh đề 3: một lần chọn đoạn Q có cận tiết lộ riêng}
Tại cùng lịch sử phân tích, thư viện và $w_a>0$ chung. Với $z_a\in[0,1]$, công thức (\ref{eq:q}) cho
\begin{equation}
\sum_ap_a=1+\alpha\left(\sum_aw_az_a-\mu\right)=1,\qquad
w_a(1-\alpha)\le p_a\le w_a(1+\alpha).
\label{eq:normalization}
\end{equation}
Do $\alpha<1$, support vẫn dương. Hai belief bất kỳ tạo hai vector điểm nhưng cùng khoảng chặn:
\begin{equation}
\boxed{\frac{p_a(b)}{p_a(b')}\le\frac{1+\alpha}{1-\alpha}=1+\varepsilon_Q\le e^{\varepsilon_Q}.}
\label{eq:qbound}
\end{equation}
Nếu belief còn ngẫu nhiên tại cùng public history, lấy kỳ vọng điểm có điều kiện vẫn nằm trong $[0,1]$;
phép centered tuyến tính nên cận vẫn giữ. Không cần chứng minh belief là posterior chính xác.

\textbf{Liên hệ mã.} Với $w_a=m_a/T$, $z_a=F_a/S$, $\alpha=N/D$, trọng số integer là
\[
W_a=m_a\left[DTS+N\left(TF_a-\sum_cm_cF_c\right)\right],\qquad\sum_aW_a=DT^2S.
\]
Một integer đều trong tổng công khai này cho đúng PMF hữu hạn. Kết luận phụ thuộc bit đều mới và
ledger đáng tin cậy; PCG dùng trong benchmark không phải chứng chỉ entropy lý thuyết thông tin.

\subsection{Mệnh đề 4: hợp thành các đoạn cho cả phiên}
\begin{equation}
\varepsilon_{Q,j}=\frac{\Gamma_Q}{j(j+1)},\qquad
\sum_{j=1}^n\varepsilon_{Q,j}=\Gamma_Q\frac{n}{n+1}\le\Gamma_Q.
\label{eq:allowance}
\end{equation}
Mã dùng tham số hữu tỷ biểu diễn được không lớn hơn allowance. Nhân các cận tại cùng augmented history:
\begin{equation}
\boxed{\Pr_X(O_n\in B)\le R_n\Pr_{X'}(O_n\in B),\qquad
R_n=\prod_{j=1}^n(1+\varepsilon_{Q,j})\le e^{\Gamma_Q}.}
\label{eq:whole}
\end{equation}
Phản hồi provider có cùng conditional law, cache suy từ replies; các khung còn lại và retry dùng bản tin đã commit.
Chứng minh trên lịch sử có cả lane/đoạn rồi marginal hóa xuống wire, không nhầm lane trùng tọa độ.
Không cần IID; không được reset bộ đếm trong cùng phiên. Với $\Gamma_Q=1$, cận hữu tỷ toàn thời lượng
nhỏ hơn $\RatioUpper$. Adaptation giảm dần khi allowance nhỏ; không cam kết utility cao cho mọi chuyến dài.

\subsection{Định lý toàn transcript và ý nghĩa suy luận}
Khi downstream chỉ dùng $E$ và thông tin công khai, kernel hậu xử lý $K$ chung cho hai traces:
\begin{equation}
\Pr_X(O\in B)=\int K(B\mid e,C)\,d\Pr_X(E=e)
\le e^{C_sD_\infty}\Pr_{X'}(O\in B).
\label{eq:postprocessing}
\end{equation}
Kết hợp hai cận cho \emph{cùng} output event:
\begin{equation}
\boxed{\Pr_X(O\in B)\le e^{\min\{C_sD_\infty,\log R_n\}}\Pr_{X'}(O\in B).}
\label{eq:minimum}
\end{equation}
Cận Q không cần upstream Geo-I đúng; nhánh không gian và minimum cần kernel Geo-I lý tưởng cùng full restore.
Bayes cho odds sau bằng odds trước nhân likelihood ratio. Với một target event có prior $\pi\in(0,1)$:
\begin{equation}
\frac{\pi}{\pi+R_n(1-\pi)}\le\Pr(Y\in A\mid O,C)
\le\frac{R_n\pi}{1-\pi+R_n\pi}.
\label{eq:bayes}
\end{equation}
Prior 50/50 cho posterior từng transcript không quá khoảng $\PosteriorUpper\%$ ở cận toàn thời lượng.
Đây không phải attack accuracy trung bình. Identity người/xe và tương lai chỉ nhận hệ quả khi target đi qua
kênh private input đã mô hình hóa; prior mạnh, ID trực tiếp và nhiều phiên liên kết vẫn cần phân tích riêng.

% PAGE 06
\clearpage
\section{Utility: bảo đảm gì và ở phạm vi nào?}
\subsection{Ưu tiên điểm hữu ích không đồng nghĩa bảo đảm Recall thật}
Tại local state cố định, với kernel (\ref{eq:q}):
\begin{equation}
\mathbb E_p[z]=\mathbb E_w[z]+\alpha\operatorname{Var}_w(z)\ge\mathbb E_w[z],\qquad
\mathbb E_p[r]-\mathbb E_w[r]=\alpha\operatorname{Cov}_w(z,r).
\label{eq:gain}
\end{equation}
$r_a$ là actual utility có điều kiện của action. Vì vậy gain của điểm chuẩn hóa luôn không âm, còn gain
actual utility cần điểm proxy tương quan đúng với chất lượng thật. Nếu $\max_a|r_a-v_a|\le\eta$ thì
\[
\Delta_r\ge\alpha\left[\operatorname{Cov}_w(z,v)-\eta\mathbb E_w|z-\mu|\right].
\]
Khi chưa chứng nhận $\eta$, không đặt bằng 0. Nếu private scores còn ngẫu nhiên, lấy kỳ vọng các
covariance có điều kiện; không thay bằng covariance của các vector đã lấy trung bình.

\subsection{Từ phản hồi thực nhận đến đáp án tại GPS thật}
Gọi $T_t$ là reference top-k theo GPS, nhu cầu và cùng total order; $U_t$ là cache metadata hợp lệ.
Lọc và sắp xếp cùng quy tắc cho
\begin{equation}
\operatorname{Recall@k}(U_t)=\frac{|T_t\cap U_t|}{|T_t|},\qquad
U_t\supseteq U_0\Longrightarrow\operatorname{Recall@k}(U_t;x_t,\psi_t)
\ge\operatorname{Recall@k}(U_0;x_t,\psi_t).
\label{eq:cache}
\end{equation}
Reference rỗng là N/A. Bất đẳng thức so \emph{cùng GPS/request}, không khẳng định Recall tăng theo thời gian
khi người dùng di chuyển. Không thể xếp hạng local để khôi phục POI chưa được lấy về. Expiry/status động
có thể phá inclusion của cache hợp lệ.

\subsection{Cận utility cả phiên so với cùng public-base planner}
Gọi $W$ là law của cùng planner khi $\Gamma_Q=0$; $P$ là law có private adaptation.
Tại cùng augmented history, $\chi^2(p_j\Vert w_j)\le\alpha_j^2/4$. Tower property và Cauchy--Schwarz cho
\begin{equation}
\TV(P_n,W_n)\le\delta_n=\frac12\sqrt{\prod_{j\le n}(1+\alpha_j^2/4)-1},\qquad
|\mathbb E_Ph-\mathbb E_Wh|\le\delta_n
\label{eq:utilitypath}
\end{equation}
với mọi $h\in[0,1]$, kể cả Recall trung bình cả phiên hoặc theo purpose. Không cần độc lập giữa các bước;
wire projection không tăng cận. Với $\Gamma_Q=1$, $\delta_\infty<0{,}0553$: tối đa 5,53 pp so với \emph{cùng}
public-base rule, không phải cận so với service-cover hay paper bất kỳ.

\textbf{Public floor sang actual utility.} Nếu mọi history có reference mean floor $c_t$, và actual case law
điều kiện theo history cách miền tham chiếu không quá $\tau_t$ theo TV, proxy error không quá $\nu_t$, thì
\begin{equation}
\mathbb E[r_t\mid H_t]\ge\max(0,c_t-\tau_t-\nu_t).
\label{eq:transfer}
\end{equation}
Tower property cho trung bình phiên với trọng số công khai, không cần IID. Unknown errors vẫn là unknown.
Floor dự đoán chỉ chuyển sang cache khi phản hồi đầy đủ và khớp signatures; benchmark cao không chứng nhận điều kiện này cho mọi provider.

\begin{tabularx}{\linewidth}{@{}P{3cm}P{5.4cm}Y@{}}\toprule
\textbf{Kết quả} & \textbf{Giả thiết/cơ chế bảo đảm} & \textbf{Giới hạn còn giữ}\\\midrule
Geo-I và cap & Miền chung, test/REM lý tưởng, prospective filter & Float sampler và full restore chưa chứng nhận.\\
Cận Q & Common support/masses, integer kernel, commit/ledger & Bit đều mới; clock và kênh quan sát đúng mô hình.\\
Utility cache & Metadata hợp lệ, cùng GPS/request/order & Không hứa floor ngoài miền; không áp cache hết hạn.\\
Utility thực & Bound proxy/TV có điều kiện hoặc so matched W & Chưa có calibration thực địa; không mặc định sai số bằng 0.\\\bottomrule
\end{tabularx}
S7 có conditional payload noninterference: đổi $\psi$ nhưng giữ cùng trajectory/context/randomness không đổi transcript.
Ý định tương quan tuyến có thể vẫn bị suy luận. Các cận inference trung bình theo mô hình target ngoại sinh
không chứng minh anonymity của tài khoản hay hành động tương lai do Q gây ra.

% PAGE 07
\clearpage
\section{Đánh giá thực nghiệm}
\subsection{Protocol, baseline và attacker trước khi đọc kết quả}
\textbf{Dữ liệu chính.} Native SUMO trên cùng bản đồ; 48 family, chia 20 TRAIN / 8 SELECTION / 20 TEST.
Hai chuyến đánh giá/family; TRAIN và SELECTION một draw, TEST ba draw: tổng 176 trip-draw, trong đó
120 TEST. Cửa sổ 0-600 s gồm 31 khung. Một family giữ cùng split khi có nhiều chuyến/draw.

GPS/hình học này đã được đánh giá trước. Protocol được đóng băng trước các draws bảo vệ mới,
không phải trước khi mọi hình học lần đầu được xem. Đây là \textbf{stochastic stability replication},
không phải xác nhận độc lập trên thành phố hoặc cohort hình học chưa xem.

\begin{tabularx}{\linewidth}{@{}P{3.8cm}Y@{}}\toprule
\textbf{Phương án} & \textbf{Vai trò và khác biệt}\\\midrule
Geo-I Z (K1) & Gửi trực tiếp cùng Z đã bảo vệ; comparator hậu kiểm có kiểm soát, một điểm/khung. Không cùng chi phí K5.\\
Service-cover K5 & Baseline bộ chọn cũ, cùng Geo-I, L60, clock và cache; không có cận Q mới như cấu hình đề xuất.\\
Initial-only K5 & Giữ bộ Q ban đầu; kiểm tra lợi ích tiếp tục khám phá.\\
Local explore & Thư viện khám phá lân cận; kiểm tra đóng góp của mục tiêu xa.\\
Public-base $\Gamma=0$ & Cùng public library/base, không điều chỉnh theo belief; control quan trọng cho utility và suy luận.\\
GeoI-BQ $\Gamma=1$ & Cấu hình chính có bounded belief adaptation, public novelty và long-range goals.\\\bottomrule
\end{tabularx}

\textbf{Điều kiện dịch vụ chung.} K5 dùng L60 mỗi loại/Q, nhịp 20 s, cùng policy cache và xếp hạng local.
Bốn purposes: gần nhất theo khoảng cách, thời gian đi ngắn nhất, trong bán kính 1000 m, ít vòng đường với
đích local đã biết. Primary dùng GPS native hiện tại mỗi event để ranking; không đồng nghĩa Geo-I đọc GPS mỗi 20 s.
L60 được chọn bằng public reference design, trước private draws mới, không chọn theo TEST score.

\textbf{Attacker bank.} Không phải tất cả đều là shadow kNN. Bank gồm \textbf{ExtraTrees64},
\textbf{Shadow kNN1/kNN5}, prior, centroid, median và geometry; endpoint bổ sung ước lượng theo slot có thứ tự và chuỗi.
Fit trên TRAIN, chọn attacker trên SELECTION \emph{riêng cho từng loss}, sau đó báo TEST.
Một ô Hit100/MAE có thể đến từ hai attacker đã chọn khác nhau. Các đặc trưng tọa độ và chuỗi thay theo scenario
và output contract, không dùng GPS thật của TEST làm feature cho attacker.

\begin{tabularx}{\linewidth}{@{}P{3cm}Y@{}}\toprule
\textbf{Metric} & \textbf{Đọc đúng}\\\midrule
Hit100 $\downarrow$ & Tỷ lệ ước lượng attacker nằm trong 100 m của target. Ví dụ 0,188\% là tỷ lệ thành công thực nghiệm, không là xác suất privacy lý thuyết.\\
MAE m $\uparrow$ & Sai số khoảng cách trung bình của attacker đã chọn. MAE cao không tự chứng minh chống mọi attacker.\\
Recall@5 $\uparrow$ & Tỷ lệ reference top-5 xuất hiện trong đáp án đã truy hồi và xếp hạng local; reference rỗng là N/A, đáp án rỗng khi có reference chấm 0.\\
CI và pp & Bootstrap ghép 10.000 lần theo family. Chênh lệch 1 pp là 1 điểm phần trăm, khác tăng tương đối 1\%.\\\bottomrule
\end{tabularx}

\textbf{Gate chọn cấu hình.} Cận dưới CI95\% của chênh Recall so baseline không thấp hơn $-2$ pp,
cả macro và từng purpose; mean mỗi draw cũng phải không thấp hơn $-2$ pp.
Đây là tiêu chí noninferiority trong protocol, không phải định lý cho population bất kỳ.

\source{Protocol/main results 14/10; comparator Z K1 cùng protected history. Mã và 108 nguồn tham chiếu được xác thực khi dựng bảng.}

% PAGE 08
\clearpage
\subsection{Kết quả privacy: đầy đủ sáu phương án, năm scenario}
Mỗi ô là \textbf{Hit100 (\%) $\downarrow$ / MAE (m) $\uparrow$}. Các giá trị là TEST point estimates
đã lưu, không tính thêm CI privacy trong lần dựng report này.
\begin{small}\par\smallskip\noindent\begin{tabularx}{\linewidth}{@{}P{3.6cm}XXXXX@{}}\toprule
\textbf{Phương án} & \textbf{S1} & \textbf{S2} & \textbf{S3} & \textbf{S9} & \textbf{S10} \\
\midrule
Geo-I Z (K1) & 18{,}522 / 338{,}5 & 37{,}500 / 214{,}3 & 18{,}522 / 338{,}5 & 5{,}833 / 435{,}8 & 35{,}000 / 246{,}8 \\
Service-cover K5 & 0{,}430 / 1145{,}8 & 0{,}833 / 1001{,}2 & 0{,}484 / 1000{,}1 & 0{,}000 / 953{,}7 & 0{,}000 / 1018{,}7 \\
Initial-only K5 & 0{,}000 / 2104{,}6 & 0{,}000 / 2124{,}2 & 0{,}000 / 2104{,}6 & 0{,}000 / 2122{,}1 & 0{,}000 / 2275{,}6 \\
Local explore & 0{,}215 / 2104{,}6 & 0{,}000 / 2124{,}2 & 0{,}215 / 2104{,}6 & 0{,}000 / 1986{,}4 & 0{,}000 / 2220{,}2 \\
Public-base $\Gamma=0$ & 0{,}054 / 2104{,}6 & 0{,}000 / 2124{,}2 & 0{,}054 / 2104{,}6 & 0{,}000 / 1986{,}4 & 0{,}000 / 2124{,}2 \\
\textbf{GeoI-BQ $\Gamma=1$} & 0{,}188 / 2104{,}6 & 0{,}000 / 2124{,}2 & 0{,}188 / 2104{,}6 & 0{,}000 / 1986{,}4 & 0{,}000 / 2124{,}2 \\
\bottomrule\end{tabularx}\par\smallskip
\end{small}

\textbf{Đọc kết quả chính.} GeoI-BQ giảm Hit100 S1 so với Z trực tiếp và service-cover; MAE của năm scenario
đều cao hơn hai comparator này. S9/S10 có Hit100 bằng 0 ở nhiều phương án nên riêng cột đó không phân biệt
được chất lượng; cần xem MAE, attack protocol và cận lý thuyết.

Initial-only có MAE S10 cao hơn GeoI-BQ, và Public-base có Hit100 S1 thấp hơn. Vì vậy không kết luận
GeoI-BQ thắng mọi control trên mọi chỉ số. Các lựa chọn phục vụ trade-off giữa thích nghi dịch vụ,
bảo đảm không gian và cận Q toàn phiên.

\textbf{Tại sao nhiều MAE bằng nhau?} Attacker theo MAE của các bounded-Q arms thường là prior median.
Bank chưa khai thác được thêm tín hiệu tốt hơn prior trong loss đó; các con số bằng nhau không chứng minh
output độc lập hoàn toàn với GPS, cũng không chứng minh đã tìm attacker tối ưu.

\textbf{Attacker thực sự được chọn cho hai phương án tiêu biểu:}
\begin{small}\par\smallskip\noindent\begin{tabularx}{\linewidth}{@{}P{3.7cm}P{1cm}YY@{}}\toprule
\textbf{Phương án} & \textbf{Case} & \textbf{Chọn theo Hit100} & \textbf{Chọn theo MAE} \\
\midrule
Geo-I Z (K1) & S1 & Tâm Q & Tâm Q \\
Geo-I Z (K1) & S2 & Tâm Q & Tâm Q \\
Geo-I Z (K1) & S3 & Tâm Q & Tâm Q \\
Geo-I Z (K1) & S9 & Tâm Q & Tâm Q \\
Geo-I Z (K1) & S10 & Tâm Q & Geometry + OLS3 (slot 0) \\
\textbf{GeoI-BQ $\Gamma=1$} & S1 & Tâm Q & Prior median \\
\textbf{GeoI-BQ $\Gamma=1$} & S2 & Tâm Q & Prior median \\
\textbf{GeoI-BQ $\Gamma=1$} & S3 & Tâm Q & Prior median \\
\textbf{GeoI-BQ $\Gamma=1$} & S9 & Tâm Q & Geometry + prior \\
\textbf{GeoI-BQ $\Gamma=1$} & S10 & Tâm Q & Geometry + prior \\
\bottomrule\end{tabularx}\par\smallskip
\end{small}
Các tên geometry/prior là baseline suy luận, không phải mô hình deep learning. Bảng thể hiện việc selector
có thể chọn control đơn giản thay vì learner, và tránh diễn giải sai rằng tất cả kết quả đều do một attacker kNN.

\textbf{Liên hệ với định lý.} Benchmark đo hiệu quả của bank hữu hạn trên workload đã khai báo.
Cận likelihood/posterior kiểm soát phân phối đầu ra dưới các giả thiết, không cần attacker trong bank tối ưu.
Hai loại bằng chứng bổ sung nhau; không dùng một score tốt để bỏ các giả thiết của định lý.

\claim{Số 0\% trong TEST không có nghĩa khả năng tấn công bằng 0. S2 mới chưa xác nhận đầy đủ tác vụ suy luận nơi dừng ngữ nghĩa.}
\source{Main results và bare-Z results 14/10; loss-specific attacker selection được giữ nguyên.}

% PAGE 09
\clearpage
\subsection{Kết quả utility và khoảng tin cậy}
\textbf{Cache phiên, Recall@5 (\%).} Tất cả K5 dùng L60 và cùng hợp đồng phục vụ; K1 là comparator có chi phí khác.
\begin{small}\par\smallskip\noindent\begin{tabularx}{\linewidth}{@{}P{4.2cm}XXXXX@{}}\toprule
\textbf{Phương án} & \textbf{Trung bình} & \textbf{Gần nhất} & \textbf{Nhanh nhất} & \textbf{Bán kính} & \textbf{Ít vòng đường} \\
\midrule
Geo-I Z (K1) & 99{,}700 & 99{,}753 & 99{,}753 & 99{,}625 & 99{,}668 \\
Service-cover K5 & 100{,}000 & 100{,}000 & 100{,}000 & 100{,}000 & 100{,}000 \\
Initial-only K5 & 98{,}377 & 98{,}984 & 98{,}987 & 96{,}358 & 99{,}181 \\
Local explore & 99{,}562 & 99{,}609 & 99{,}609 & 99{,}358 & 99{,}670 \\
Public-base $\Gamma=0$ & 99{,}711 & 99{,}730 & 99{,}731 & 99{,}642 & 99{,}740 \\
\textbf{GeoI-BQ $\Gamma=1$} & 99{,}595 & 99{,}638 & 99{,}641 & 99{,}388 & 99{,}712 \\
\bottomrule\end{tabularx}\par\smallskip
\end{small}
GeoI-BQ đạt \textbf{\PrimaryRecall\%}; CI95\% macro là [99,225; 99,902]\%.
So service-cover, chênh lệch nhỏ nhưng âm; không gọi đó là vượt baseline về utility.

\textbf{Độ bất định theo purpose và gate đã định trước:}
\begin{small}\par\smallskip\noindent\begin{tabularx}{\linewidth}{@{}P{3.1cm}XXXX@{}}\toprule
\textbf{Purpose} & \textbf{BQ (\%)} & \textbf{CI tuyệt đối 95\%} & \textbf{Δ so K5 (pp)} & \textbf{CI Δ (pp)} \\
\midrule
Trung bình & 99{,}595 & [99{,}225; 99{,}902] & -0{,}405 & [-0{,}775; -0{,}098] \\
Gần nhất & 99{,}638 & [99{,}269; 99{,}927] & -0{,}362 & [-0{,}731; -0{,}073] \\
Nhanh nhất & 99{,}641 & [99{,}276; 99{,}927] & -0{,}359 & [-0{,}724; -0{,}073] \\
Bán kính & 99{,}388 & [98{,}669; 99{,}903] & -0{,}612 & [-1{,}331; -0{,}097] \\
Ít vòng đường & 99{,}712 & [99{,}358; 99{,}946] & -0{,}288 & [-0{,}642; -0{,}054] \\
\bottomrule\end{tabularx}\par\smallskip
\end{small}
Cận dưới yếu nhất của CI chênh lệch là radius, khoảng $-1{,}331$ pp; vẫn trên $-2$ pp.
Cấu hình chính qua gate macro/từng purpose và kiểm tra mean từng draw trong artifact gốc.
Không suy CI này thành floor Recall cho mọi chuyến.

\textbf{Chênh lệch macro của GeoI-BQ so từng control:}
\begin{small}\par\smallskip\noindent\begin{tabularx}{\linewidth}{@{}P{6cm}XX@{}}\toprule
\textbf{Đối chứng} & \textbf{BQ trừ đối chứng (pp)} & \textbf{CI ghép 95\% (pp)} \\
\midrule
Geo-I Z (K1) & -0{,}105 & [-0{,}536; 0{,}273] \\
Service-cover K5 & -0{,}405 & [-0{,}775; -0{,}098] \\
Initial-only K5 & 1{,}218 & [0{,}302; 2{,}412] \\
Local explore & 0{,}033 & [-0{,}109; 0{,}178] \\
Public-base $\Gamma=0$ & -0{,}116 & [-0{,}253; -0{,}006] \\
\bottomrule\end{tabularx}\par\smallskip
\end{small}
\begin{itemize}
\item So initial-only, gain khoảng 1,218 pp và CI không chạm 0: tiếp tục khám phá có ích trong workload này.
\item So local explore, CI chạm 0: trial L60 chưa xác nhận một gain riêng từ long-range goals.
\item Public-base $\Gamma=0$ có Recall cao hơn khoảng 0,116 pp; CI của BQ trừ control âm.
Không có bằng chứng belief adaptation tăng actual utility trong catalogue này.
\item So bare Geo-I Z, CI chứa 0; không chứng minh thắng utility, càng không chứng minh thắng utility/cost.
\end{itemize}

\claim{Kết quả hỗ trợ một cấu hình có utility cao và qua gate đã chọn, không phải một mô hình tối ưu mọi mặt. Vai trò của Geo-I là bảo đảm không gian, không phải nguyên nhân được chứng minh của mọi retrieval gain.}
\source{Main session-cache metrics; paired-family CIs trong release readout; bare-Z controlled comparison.}

% PAGE 10
\clearpage
\subsection{Chi phí, cache và độ nhạy của kết quả}
\textbf{Chi phí trung bình mỗi phiên 600 s.} KB = 1000 bytes; query điểm là một tọa độ gửi server.
\begin{small}\par\smallskip\noindent\begin{tabularx}{\linewidth}{@{}P{4.4cm}XXX@{}}\toprule
\textbf{Phương án} & \textbf{Số query điểm} & \textbf{Request (KB)} & \textbf{Reply-index (KB)} \\
\midrule
Geo-I Z (K1) & 31 & 6{,}352 & 16{,}921 \\
Service-cover K5 & 155 & 31{,}756 & 85{,}268 \\
Initial-only K5 & 155 & 31{,}707 & 85{,}318 \\
Local explore & 155 & 31{,}748 & 82{,}788 \\
Public-base $\Gamma=0$ & 155 & 31{,}751 & 83{,}442 \\
\textbf{GeoI-BQ $\Gamma=1$} & 155 & 31{,}756 & 82{,}557 \\
\bottomrule\end{tabularx}\par\smallskip
\end{small}
Đây là byte JSON request và reply \emph{indices}; chưa gồm metadata đầy đủ, HTTP/TLS, latency, CPU, RAM đỉnh hay pin.
K5 gửi 155 query điểm so 31 của K1. Không coi hai output contracts là cùng chi phí.

\textbf{Ba cách dùng phản hồi, macro Recall (\%):}
\begin{small}\par\smallskip\noindent\begin{tabularx}{\linewidth}{@{}P{4.3cm}XXX@{}}\toprule
\textbf{Phương án} & \textbf{Current (\%)} & \textbf{Epoch 60 s (\%)} & \textbf{Session (\%)} \\
\midrule
Service-cover K5 & 99{,}831 & 99{,}872 & 100{,}000 \\
Initial-only K5 & 98{,}377 & 98{,}377 & 98{,}377 \\
Local explore & 96{,}671 & 96{,}941 & 99{,}562 \\
Public-base $\Gamma=0$ & 97{,}078 & 97{,}380 & 99{,}711 \\
\textbf{GeoI-BQ $\Gamma=1$} & 96{,}780 & 97{,}123 & 99{,}595 \\
\bottomrule\end{tabularx}\par\smallskip
\end{small}
Current chỉ dùng phản hồi hiện tại; epoch giữ trong cửa sổ 60 s; session giữ metadata tĩnh hợp lệ cả phiên.
Gate primary áp cho session cache. Chỉ dùng current, BQ đạt khoảng 96,780\% và không qua gate đó.

\textbf{Cache bao phủ catalogue 418 POI:}
\begin{small}\par\smallskip\noindent\begin{tabularx}{\linewidth}{@{}P{4.3cm}XXX@{}}\toprule
\textbf{Phương án} & \textbf{Cache đầu} & \textbf{Cache cuối} & \textbf{Recall 200-600 s (\%)} \\
\midrule
Service-cover K5 & 359{,}02 & 377{,}41 & 100{,}000 \\
Initial-only K5 & 380{,}20 & 380{,}20 & 98{,}240 \\
Local explore & 380{,}20 & 411{,}22 & 99{,}823 \\
Public-base $\Gamma=0$ & 380{,}12 & 416{,}08 & 99{,}986 \\
\textbf{GeoI-BQ $\Gamma=1$} & 380{,}20 & 415{,}88 & 99{,}881 \\
\bottomrule\end{tabularx}\par\smallskip
\end{small}
Cache BQ trung bình từ 380,20 lên 415,88/418 POI. Near-full cache giải thích utility rất cao, đồng thời hạn chế
khả năng ngoại suy sang catalogue lớn, trạng thái thay đổi hoặc workload khác. Bulk-catalogue local ranking là control
cần thảo luận khi API cho phép tải toàn bộ, chưa có thí nghiệm matched mới chứng minh BQ tốt hơn contract bulk.

\textbf{GPS dùng ranking chỉ cập nhật mỗi 60 s, giữ fix cuối, Recall (\%):}
\begin{small}\par\smallskip\noindent\begin{tabularx}{\linewidth}{@{}P{4.3cm}XXXXX@{}}\toprule
\textbf{Phương án} & \textbf{Trung bình} & \textbf{Gần nhất} & \textbf{Nhanh nhất} & \textbf{Bán kính} & \textbf{Ít vòng đường} \\
\midrule
Service-cover K5 & 95{,}625 & 97{,}126 & 97{,}129 & 89{,}902 & 98{,}343 \\
Public-base $\Gamma=0$ & 95{,}432 & 96{,}938 & 96{,}941 & 89{,}706 & 98{,}143 \\
\textbf{GeoI-BQ $\Gamma=1$} & 95{,}348 & 96{,}873 & 96{,}879 & 89{,}514 & 98{,}124 \\
\bottomrule\end{tabularx}\par\smallskip
\end{small}
Giữ nguyên Q và privacy budget; chỉ đổi vị trí dùng xếp hạng local. Macro BQ còn 95,348\%, radius 89,514\%.
Không chuyển con số primary 99,595\% sang bối cảnh GPS local cũ này. Đây chưa là dữ liệu sensor thực địa được hiệu chuẩn.

\claim{Utility cao phụ thuộc cache hợp lệ và vị trí local đủ mới. Chi phí đầy đủ và khả năng triển khai điện thoại vẫn chưa được đo.}

% PAGE 11
\clearpage
\subsection{Độ sâu phản hồi: thiết kế công khai và ablation}
L60 được chọn từ bảng tham chiếu công khai gồm 36 trạng thái, bốn purposes và chín destination prototypes,
không đọc GPS/benchmark score để chọn L. Chọn L nhỏ nhất trong frontier đã thử mà public beam có goal đạt mean floor 99\%:
\begin{small}\par\smallskip\noindent\begin{tabularx}{\linewidth}{@{}XXX@{}}\toprule
\textbf{L / loại / Q} & \textbf{Mean miền công khai (\%)} & \textbf{Goal đạt floor 99\%} \\
\midrule
30 & 93{,}727 & 0 \\
40 & 96{,}651 & 0 \\
50 & 98{,}401 & 0 \\
60 & 99{,}351 & 24 \\
\bottomrule\end{tabularx}\par\smallskip
\end{small}
Đây là minimum trong beam/frontier đã khai báo, không phải optimum toàn mạng và không phải QoS 99\% cho mọi GPS.
Initial pool chỉ nhận floor khi các phản hồi thực nhận đầy đủ và khớp signatures.

\textbf{Giữ nguyên Q được lập kế hoạch ở L60, chỉ cắt phản hồi còn L30:}
\begin{small}\par\smallskip\noindent\begin{tabularx}{\linewidth}{@{}P{3.3cm}XXX@{}}\toprule
\textbf{Purpose} & \textbf{L30 cùng Q (\%)} & \textbf{L60 trừ L30 (pp)} & \textbf{CI Δ (pp)} \\
\midrule
Trung bình & 94{,}170 & 5{,}425 & [3{,}486; 7{,}527] \\
Gần nhất & 96{,}105 & 3{,}534 & [2{,}371; 4{,}726] \\
Nhanh nhất & 96{,}101 & 3{,}540 & [2{,}375; 4{,}740] \\
Bán kính & 89{,}015 & 10{,}374 & [5{,}768; 15{,}706] \\
Ít vòng đường & 95{,}459 & 4{,}253 & [2{,}946; 5{,}520] \\
\bottomrule\end{tabularx}\par\smallskip
\end{small}
Macro cùng-Q L30 là 94,170\%, thấp hơn L60 khoảng 5,425 pp; radius bị ảnh hưởng mạnh nhất.
Đây là counterfactual theo cùng wire, \emph{không phải} chạy lại một policy L30 có cache/thư viện thích nghi khác.
Không dùng bảng này để tuyên bố mọi gain là do thuật toán mới hoặc do ngân sách privacy.

\section{Thảo luận: đóng góp và các điều kiện chưa đóng}
\begin{tabularx}{\linewidth}{@{}P{3cm}Y@{}}\toprule
\textbf{Phạm vi} & \textbf{Đánh giá hiện tại}\\\midrule
Đã có lập luận & REM/noisy test lý tưởng; cap và composition; bounded Q kernel; cận toàn transcript và posterior; utility/cache/path theo điều kiện.\\
Chưa chứng nhận mã & Geo-I float có thể có support khác lý tưởng. Timestamp/config không xác thực full upstream restore; reset sai có thể làm sai accounting Geo-I. Cận Q có bộ giả thiết riêng.\\
Giới hạn đánh giá & Cùng bản đồ, geometry đã xem, catalogue nhỏ/tĩnh, exact local GPS ở primary; bank hữu hạn và chưa có field deployment.\\
S4--S8 & Inference có điều kiện và payload noninterference không đồng nghĩa anonymity của người/xe hay bảo vệ mọi ý định/người đồng hành.\\\bottomrule
\end{tabularx}

\section{Kết luận theo ba câu hỏi nghiên cứu}
\textbf{RQ1.} Có chuỗi lập luận từ thao tác GPS đến toàn transcript, với cận không gian Geo-I và cận Q riêng.
Không cần chứng minh ``không thể đảo Q về Z''; điều cần chặn là bằng chứng quan sát giúp phân biệt các traces.

\textbf{RQ2.} Chọn Q ưu tiên điểm proxy; cache và local ranking có bảo đảm rõ phạm vi.
Cận actual utility cần matched base hoặc sai số transfer có điều kiện, không được thay bằng benchmark mean.

\textbf{RQ3.} Cấu hình chính đạt utility cao và gate $-2$ pp trong protocol, với attack scores vị trí thuận lợi hơn
bare Z/service-cover ở các chỉ số đã nêu. Nó không thắng mọi control; chi phí K5 lớn hơn K1 và $\Gamma=0$ vẫn có Recall tốt hơn.

Đóng góp nằm ở đặc tả, thiết kế lựa chọn đoạn có cận và đánh giá tái lập; Geo-I, REM, reuse, composition và
postprocessing là nền tảng kế thừa. Bản luận văn tiếp theo cần tích hợp nhất quán định lý--giả thiết--mã,
không thêm cơ chế chỉ để tối ưu điểm TEST. Phụ lục giữ evidence lịch sử và pilot riêng để review đầy đủ.

% PAGE 12
\clearpage
\appendix
\section{Đối chứng paper và metric lịch sử: không ghép vào trial chính}
\textbf{Paper-v2, K5, 12 chuyến test, ba seed SUMO.} Mỗi ô: Hit100\% / MAE m / Recall@5\%.
Dấu * chỉ local adaptation; không tuyên bố tái lập SOTA đầy đủ.
\begin{small}\par\smallskip\noindent\begin{tabularx}{\linewidth}{@{}P{4cm}XXX@{}}\toprule
\textbf{Phương pháp} & \textbf{S1} & \textbf{S2} & \textbf{S3} \\
\midrule
GPS thật & 100{,}00 / 13 / 100{,}00 & 100{,}00 / 39 / 100{,}00 & 96{,}43 / 23 / 100{,}00 \\
DLS* & 58{,}33 / 1318 / 100{,}00 & 100{,}00 / 39 / 100{,}00 & 31{,}21 / 1688 / 100{,}00 \\
TransProtect* & 83{,}33 / 84 / 97{,}78 & 66{,}67 / 88 / 97{,}19 & 69{,}28 / 99 / 97{,}93 \\
Semantic* & 100{,}00 / 11 / 100{,}00 & 100{,}00 / 42 / 100{,}00 & 40{,}58 / 257 / 100{,}00 \\
BR--Geo-I lịch sử & 33{,}33 / 299 / 95{,}83 & 22{,}22 / 247 / 97{,}01 & 7{,}44 / 776 / 87{,}67 \\
\bottomrule\end{tabularx}\par\smallskip
\end{small}
BR ở đây là \texttt{br\_private} lịch sử, không phải GeoI-BQ v1, không phải L60 trial.
Các output contracts và ngân sách khác nhau. DLS có MAE S1/S3 lớn hơn BR; BR có Hit100 thấp hơn nhưng Recall S3 thấp hơn.
Không lấy các cột từ cohort này để tính gain so với bảng mới.

\subsection{Metric gốc khả dụng trên cùng bản thích nghi}
Readout native-metric chỉ tái sử dụng execution đã lưu, không fit lại. Bảng dưới lấy S3 cùng paper-v2:
\begin{small}\par\smallskip\noindent\begin{tabularx}{\linewidth}{@{}P{4.4cm}XXXX@{}}\toprule
\textbf{Phương pháp} & \textbf{Hoàn tất / thử} & \textbf{EIE (m)} & \textbf{Dịch chuyển (m)} & \textbf{Sinh / event (ms)} \\
\midrule
GPS thật & 12 / 12 & 23{,}49 & 23{,}49 & 0{,}01 \\
DLS* & 12 / 12 & 1687{,}67 & -- & 1{,}10 \\
TransProtect* & 12 / 12 & 99{,}03 & 97{,}66 & 0{,}35 \\
Semantic* & 12 / 12 & 257{,}04 & -- & 1{,}28 \\
BR--Geo-I lịch sử & 12 / 12 & 776{,}32 & -- & 6{,}37 \\
AnotherMe* (offline) & 11 / 12 & 2355{,}89 & 1858{,}83 & 1{,}62 \\
\bottomrule\end{tabularx}\par\smallskip
\end{small}
\textbf{EIE của point estimator} chính là expected distance error, tương ứng MAE đã có, không tạo ra một bằng chứng privacy mới.
Single-release displacement đo khoảng cách output thay thế theo contract khả dụng; không định nghĩa tùy ý cho tập dummy.
Generation time là thời gian simulator đã ghi, không phải latency server hoặc điện thoại.

AnotherMe* yêu cầu complete route; S1/S2 không áp dụng trong cohort này. S3 hoàn tất 11/12, một failure không
được tính là privacy thành công. Các số AnotherMe trong bảng chỉ có điều kiện trên execution hoàn tất.
Những metric đòi posterior attacker, mô hình learned/calibrated gốc hoặc thông tin không được xuất phải để N/A,
không chế biến từ số dummy hoặc thay metric để có một bảng đủ ô.

\subsection{Cách dùng bằng chứng này trong luận văn}
\begin{enumerate}
\item Chương related works trình bày output contract, mục tiêu bảo vệ và metric gốc; giải thích phần nào tương thích với protocol chung.
\item Chương thực nghiệm giữ bảng lịch sử này như controlled adaptations, còn bảng GeoI-BQ v1 dùng matched internal controls trên cohort mới hơn.
\item Không tuyên bố cấu hình hiện tại vượt trực tiếp tất cả paper khi chưa chạy chúng trên cùng protocol mới với fidelity và chi phí tương ứng.
\end{enumerate}

\textbf{Khoảng trống thực nghiệm còn rõ:} chưa có một lần xác nhận mới so toàn bộ paper adaptations với GeoI-BQ v1
trên geometry chưa xem. Luận văn có thể báo bằng chứng hiện có và giới hạn này; không gọi benchmark lịch sử là
head-to-head confirmation của phương pháp cuối.

\source{Paper benchmark results; native metrics readout/verification 05/10. Số liệu là evidence lịch sử, không thay dữ liệu hay đánh giá lại.}

% PAGE 13
\clearpage
\section{Pilot continual Geo-I: evidence bổ sung, chưa thay cấu hình chính}
Một draw mới/chuyến trên 48 family; 10 phút đầu tái sử dụng native GPS, kéo dài bằng đường đi tổng hợp có hướng
đến 30 phút. TRAIN/SELECTION/TEST giữ 20/8/20 family. Phần sau không phải GPS thực địa hoặc native SUMO mới.
Ba arms cùng K5/L60, clock và cache; capped giữ $C_s=0{,}23$, continual giữ $u=0{,}01$ cho mỗi thao tác nhưng bỏ cap Geo-I.

\textbf{Macro Recall với cache phiên theo cửa sổ:}
\begin{small}\par\smallskip\noindent\begin{tabularx}{\linewidth}{@{}P{4.2cm}XXX@{}}\toprule
\textbf{Pilot} & \textbf{0-600 s (\%)} & \textbf{0-1800 s (\%)} & \textbf{720-1800 s (\%)} \\
\midrule
Giữ cap, $\Gamma=1$ & 99{,}633 & 99{,}873 & 99{,}997 \\
Continual, $\Gamma=1$ & 99{,}633 & 99{,}873 & 99{,}997 \\
Continual, $\Gamma=0$ & 99{,}692 & 99{,}891 & 99{,}994 \\
\bottomrule\end{tabularx}\par\smallskip
\end{small}
\textbf{Đầy đủ bốn purposes cho toàn 30 phút:}
\begin{small}\par\smallskip\noindent\begin{tabularx}{\linewidth}{@{}P{4.2cm}XXXXX@{}}\toprule
\textbf{Pilot, cache phiên} & \textbf{Trung bình} & \textbf{Gần nhất} & \textbf{Nhanh nhất} & \textbf{Bán kính} & \textbf{Ít vòng đường} \\
\midrule
Giữ cap, $\Gamma=1$ & 99{,}873 & 99{,}905 & 99{,}905 & 99{,}858 & 99{,}825 \\
Continual, $\Gamma=1$ & 99{,}873 & 99{,}905 & 99{,}905 & 99{,}858 & 99{,}825 \\
Continual, $\Gamma=0$ & 99{,}891 & 99{,}896 & 99{,}896 & 99{,}890 & 99{,}882 \\
\bottomrule\end{tabularx}\par\smallskip
\end{small}
\textbf{Độ bất định macro so capped:}
\begin{small}\par\smallskip\noindent\begin{tabularx}{\linewidth}{@{}P{4cm}XX@{}}\toprule
\textbf{Pilot, toàn 30 phút} & \textbf{Δ so cap (pp)} & \textbf{CI 95\% (pp)} \\
\midrule
Giữ cap, $\Gamma=1$ & 0{,}000 & [0{,}000; 0{,}000] \\
Continual, $\Gamma=1$ & 0{,}000 & [0{,}000; 0{,}000] \\
Continual, $\Gamma=0$ & 0{,}018 & [-0{,}018; 0{,}063] \\
\bottomrule\end{tabularx}\par\smallskip
\end{small}
\textbf{Hit100\% / MAE m, toàn 30 phút:}
\begin{small}\par\smallskip\noindent\begin{tabularx}{\linewidth}{@{}P{3.6cm}XXXXX@{}}\toprule
\textbf{Pilot} & \textbf{S1} & \textbf{S2} & \textbf{S3} & \textbf{S9} & \textbf{S10} \\
\midrule
Giữ cap, $\Gamma=1$ & 1{,}264 / 1712{,}0 & 2{,}500 / 705{,}1 & 1{,}264 / 1610{,}1 & 0{,}000 / 1986{,}4 & 0{,}000 / 718{,}9 \\
Continual, $\Gamma=1$ & 1{,}264 / 1712{,}0 & 2{,}500 / 705{,}1 & 1{,}264 / 1611{,}3 & 0{,}000 / 1986{,}4 & 0{,}000 / 718{,}9 \\
Continual, $\Gamma=0$ & 1{,}264 / 1712{,}0 & 0{,}000 / 705{,}1 & 1{,}264 / 1712{,}0 & 0{,}000 / 1986{,}4 & 0{,}000 / 1024{,}1 \\
\bottomrule\end{tabularx}\par\smallskip
\end{small}
\textbf{Chi phí cùng định nghĩa reply-index JSON:}
\begin{small}\par\smallskip\noindent\begin{tabularx}{\linewidth}{@{}P{4cm}XXX@{}}\toprule
\textbf{Pilot} & \textbf{Query điểm} & \textbf{Request KB} & \textbf{Reply-index KB} \\
\midrule
Giữ cap, $\Gamma=1$ & 455 & 93{,}76 & 228{,}69 \\
Continual, $\Gamma=1$ & 455 & 93{,}76 & 228{,}69 \\
Continual, $\Gamma=0$ & 455 & 93{,}77 & 228{,}73 \\
\bottomrule\end{tabularx}\par\smallskip
\end{small}

Số private reads trung bình tăng từ 13,20 lên 31,00; chi phí nhánh ẩn trung bình từ 0,2245 lên 0,5523 m$^{-1}$.
Cận prefix continual với $N=31$ là $(2N-1)u=0{,}61$ m$^{-1}$, không còn cap 0,23 cho toàn thời lượng.
Chi phí nhánh thực tế là diagnostic local, không phải certificate công khai thay cho cận prefix.

Pilot chưa cho thấy gain Recall từ việc tiếp tục đọc GPS; cache tĩnh gần đầy có thể che lợi ích thích nghi.
Không gán 99,873\% của pilot cho cohort chính hoặc diễn giải là xác nhận arbitrary-duration QoS.
Cận Q toàn phiên vẫn theo bộ giả thiết riêng; float/PCG chưa là chứng chỉ pure Geo-I.
Validation đã lưu đạt 26.208 Q-prefix checks và 129.600 directed transitions; không tự động promote.

\source{Continual pilot results/protocol/validation 17/10. Tất cả bảng của report được xuất từ evidence đã lưu; không benchmark lại.}
\end{document}
